% =====================================================================
%  Ký hiệu dùng chung cho cả báo cáo (P1 quản lý).
%  Quy ước: vector chữ thường in đậm, ma trận chữ hoa in đậm,
%  vô hướng chữ thường nghiêng. Muốn thêm macro → báo P1, đừng tự định nghĩa
%  trong file section để tránh trùng tên.
% =====================================================================

% ---- Tập số, không gian ----
\newcommand{\R}{\mathbb{R}}
\newcommand{\N}{\mathbb{N}}
\newcommand{\X}{\mathcal{X}}               % không gian đầu vào
\newcommand{\Hk}{\mathcal{H}_k}            % RKHS của hàm nhân k

% ---- Vector, ma trận ----
\newcommand{\vect}[1]{\bm{#1}}
\newcommand{\vw}{\vect{w}}
\newcommand{\vx}{\vect{x}}
\newcommand{\vz}{\vect{z}}
\newcommand{\vy}{\vect{y}}
\newcommand{\vu}{\vect{u}}
\newcommand{\vc}{\vect{c}}
\newcommand{\valpha}{\vect{\alpha}}
\newcommand{\vxi}{\vect{\xi}}
\newcommand{\vmu}{\vect{\mu}}
\newcommand{\ve}{\vect{e}}                 % vector toàn số 1
\newcommand{\vzero}{\vect{0}}
\newcommand{\mK}{\mathbf{K}}               % ma trận Gram, K_ij = k(x_i, x_j)
\newcommand{\mQ}{\mathbf{Q}}               % Q_ij = y_i y_j K_ij
\newcommand{\mI}{\mathbf{I}}
\newcommand{\T}{^{\top}}                   % chuyển vị: \vw\T\vx

% ---- Chuẩn, tích vô hướng ----
\newcommand{\norm}[1]{\left\lVert #1 \right\rVert}
\newcommand{\abs}[1]{\left\lvert #1 \right\rvert}
\newcommand{\inner}[2]{\left\langle #1, #2 \right\rangle}
\newcommand{\Hnorm}[1]{\norm{#1}_{\Hk}}
\newcommand{\Hinner}[2]{\inner{#1}{#2}_{\Hk}}

% ---- Toán tử ----
\DeclareMathOperator*{\argmin}{arg\,min}
\DeclareMathOperator*{\argmax}{arg\,max}
\DeclareMathOperator{\sign}{sign}
\DeclareMathOperator{\tr}{tr}
\DeclareMathOperator{\spn}{span}
\newcommand{\ind}[1]{\mathbf{1}\!\left[#1\right]}   % hàm chỉ thị
\newcommand{\hinge}[1]{\max\{0,\, 1 - #1\}}         % mất mát hinge

% ---- SVM, SMO ----
\newcommand{\Iup}{I_{\mathrm{up}}}
\newcommand{\Ilow}{I_{\mathrm{low}}}
\newcommand{\SV}{\mathcal{S}}               % tập chỉ số vector hỗ trợ {i : α_i > 0}
\newcommand{\Dual}{D}                       % hàm mục tiêu đối ngẫu (dạng max)
\newcommand{\nsv}{n_{\mathrm{SV}}}          % số vector hỗ trợ
